\documentclass[11pt]{article}
\usepackage{hyperref}
\usepackage{enumitem}
\usepackage{amsfonts} 
\setlength{\oddsidemargin}{-0.1in}
\setlength{\evensidemargin}{\oddsidemargin}
\setlength{\textwidth}{7.0in}
\setlength{\textheight}{9.3in}
\setlength{\topmargin}{-0.7in}
%\input{macros.tex}
\begin{document}



\fbox{
\vbox{
\begin{flushleft}
Ann,  Bob, Charlie    (\emph{replace with your names and change the date})\\  % authors' names
COSC 417 \\  %class
11/6/2021\\  % date
\end{flushleft}
\center{\Large{\textbf{Assignment 7}}}
} % end vbox
} % end fbox
\vline

%\fbox{\vbox{
\textbf{Instructions.}
\begin{enumerate}

\item  Due date and time: As announced on Blackboard.

\item %Work in teams as in the previous assignments (hand in one assignment with all the names on it).
Submit on Blackboard.  This homework will not be graded, I will post solutions (see the syllabus for the policy regarding assignments).




\item For editing your homework, I recommend that you use Latex on Overleaf, see the template files posted on Blackboard. 
 
         The template has an  example showing how to do the state diagram  of a finite automaton. If you have trouble with making figures in Latex, you can hand-draw the state diagrams (or other figures) and insert a digital photo in the .tex file (see the template how to do this).
  
\item If a problem has more questions, write down your answers in the same order as the order of questions. In principle, this should help you.

 \item   To append in the  latex file a .jpg file (for a photo; for example, in case you draw a picture by hand and take a photo of it with your phone camera), use 


\begin{verbatim}
\includegraphics[angle=270,origin=c,width=\linewidth]{file.jpg}
\end{verbatim}
The parameter angle=270 is for rotating the photo, and you may have to change 270 to whatever angle works for your photo.
		   \vspace{0.5cm}


\end{enumerate}

\newpage




\begin{enumerate}[label=\textbf{Exercise \arabic{*}}]

\item $A$ and $B$ are two languages, and $A \le_m B$ ( $\le_m$ is the mapping reduction defined in Sipser, Def 5.2, page 235, and also in my Notes 11). Show that if $B$ is Turing-recognizable, then $A$ is also Turing recognizable.

\medskip



\textbf{Answer:}. 

\item Let $B =\{\langle<M\rangle \mid \textrm{$M$ is a Turing machine that accepts 0 , and that does  not accept 1}\}$. Show that $HALT \le_m B$.
\medskip 

\textbf{Answer:}

\item Give an example of an undecidable language $B$ with $B \le_m \overline{B}$.
\medskip

\textbf{Answer}
\end{enumerate}




 \end{document}
 